\documentclass[12pt,a4paper]{article}
\usepackage[margin=2.5cm]{geometry}
\usepackage{graphicx}
\usepackage{xcolor}
\usepackage{fancyhdr}
\setlength{\headheight}{15pt}
\addtolength{\topmargin}{-3pt}
\usepackage{array}
\usepackage{longtable}
\usepackage{hyperref}
\usepackage{amssymb}

% Color definitions. Identicas a las de la ERS y el plan de alcance: los tres
% documentos son piezas del mismo entregable y deben verse como tales.
\definecolor{mainRed}{HTML}{EC3013}
\definecolor{mainGray}{HTML}{6b6866}
\definecolor{mainWhite}{HTML}{F3F2F2}

% Header and footer
\pagestyle{fancy}
\fancyhf{}
\fancyhead[C]{\textcolor{mainGray}{Declaración de Diseño del Frontend — MSO}}
\fancyfoot[C]{\thepage}
\renewcommand{\headrulewidth}{0.5pt}
\renewcommand{\headrule}{\color{mainRed}\hrule}

% Saltos de titulo mas cortos que los de article, por el mismo motivo que en el
% Plan de Gestion del Alcance: este documento acumula subsecciones cortas y el
% espaciado por omision las estira sin anadir informacion.
\makeatletter
\renewcommand\section{\@startsection{section}{1}{\z@}%
  {-2.2ex \@plus -1ex \@minus -.2ex}{1.0ex \@plus.2ex}%
  {\normalfont\Large\bfseries}}
\renewcommand\subsection{\@startsection{subsection}{2}{\z@}%
  {-1.7ex \@plus -.8ex \@minus -.2ex}{0.5ex \@plus.2ex}%
  {\normalfont\large\bfseries}}
\def\@listi{\leftmargin\leftmargini \parsep 0\p@ \topsep 4\p@ \itemsep 2\p@}
\let\@listI\@listi
\makeatother

% Title formatting
\newcommand{\sectiontitle}[1]{%
  \section*{#1}
  \addcontentsline{toc}{section}{#1}
}

\newcommand{\subsectiontitle}[1]{%
  \subsection*{#1}
  \addcontentsline{toc}{subsection}{#1}
}

% Rutas y codigos. \texttt es indivisible y una ruta larga se sale de la caja.
\urlstyle{tt}
\makeatletter
\g@addto@macro\UrlBreaks{\do\-}
\makeatother
\newcommand{\ruta}[1]{\nolinkurl{#1}}

% Marcas de estado, con el mismo significado y color que en el apartado 3.2 de
% la ERS, de donde proceden.
\newcommand{\estadoImplementado}{\textcolor{mainRed}{\textbf{[IMPLEMENTADO]}}}
\newcommand{\estadoPrevisto}{\textcolor{mainGray}{\textbf{[PREVISTO]}}}

% Marca de decision. Este documento decide, y conviene que se vea cual es la
% frase que decide y cual la que explica.
\newcommand{\decision}[1]{\textcolor{mainRed}{\textbf{#1}}}

\setlength{\parindent}{0pt}
\setlength{\parskip}{6pt}

\begin{document}

\begin{titlepage}
  \centering
  \vspace*{2cm}

  {\Large\textbf{UNIVERSIDAD DISTRITAL}\\
   \textbf{FRANCISCO JOSÉ DE CALDAS}}\par

  \vspace{3cm}

  {\Huge\textbf{Declaración de Diseño del Frontend}}\par

  \vspace{2cm}

  {\Large\textbf{MalphasOS}}\\
  {\large Gestión de Clientes}\par

  \vspace{3cm}

  {\large Sean Sebastián Bolívar Calderón - 20231020135}\par

  \vfill

  {\large 13 de septiembre de 2026}
\end{titlepage}

\tableofcontents
\newpage

% =====================================================================
\sectiontitle{Nota preliminar}

\textbf{Qué es este documento.} La declaración de las decisiones de diseño del
frontend de MalphasOS: plataforma, arquitectura, sistema visual, accesibilidad,
comportamiento sin conexión y estrategia de pruebas. Cada decisión va con su
porqué y con su coste, porque una decisión sin coste escrito se relee más tarde
como si no lo hubiera tenido.

\textbf{Qué no es.} No es un manual de estilo ni una guía de componentes: eso se
escribirá cuando los componentes existan. Tampoco es una lista de pantallas. Y
no sustituye a la ERS: \decision{este documento no introduce ningún requisito
nuevo}, sólo decide cómo se satisfacen los que ya están escritos.

\textbf{Por qué existe y por qué ahora.} El backend está terminado en sus cinco
módulos y no hay una sola línea de frontend. Esa asimetría tiene un efecto
medible: de los 31 requisitos funcionales de la ERS hay 11 implementados, y
\textbf{cuatro de los que faltan en órdenes de trabajo están abiertos únicamente
porque no existe el formulario}, no porque falte nada en el servidor. Seis
requisitos no funcionales enteros —RNF-12 a RNF-17— no tienen relación alguna
con el backend. Un backend completo con cero frontend cumple cero de ellos.

\textbf{De dónde salen las restricciones.} De la ERS, no de la preferencia de
quien escribe. El apartado 2.1 exige que el sistema sea accesible desde el
teléfono del ingeniero \textbf{sin instalar aplicaciones nativas}; el 3.1.2 lista
la pantalla táctil entre las interfaces de hardware; el 3.1.1 recoge los seis
requisitos de interfaz. Donde este documento decide algo que la ERS no dice, se
señala explícitamente como decisión propia.

% =====================================================================
\sectiontitle{1. Alcance}

\subsectiontitle{1.1. Qué cubre}

La \textbf{aplicación operativa} de MalphasOS: la interfaz con la que los
administradores gestionan clientes, equipos y órdenes de trabajo, y con la que
los ingenieros consultan y ejecutan su trabajo en campo. Cubre las decisiones
transversales que afectan a todas sus pantallas.

\subsectiontitle{1.2. Qué queda fuera, y por qué}

\textbf{La \textit{landing page} de la empresa queda fuera de esta
declaración.} El cronograma del documento de constitución la sitúa en la semana
1, de modo que su ausencia aquí podría leerse como un olvido. No lo es:

\begin{itemize}
  \item Es \textbf{otro producto}. Es pública, no tiene autenticación, no
        consume el API y su objetivo es comercial, no operativo.
  \item \textbf{No avanza ninguno de los seis requisitos de interfaz.} RNF-12 a
        RNF-17 hablan de la aplicación que un ingeniero usa, no de una página
        de presentación.
  \item \textbf{Pide un stack distinto.} Una página pública necesita
        indexación por buscadores, y una aplicación de página única es
        precisamente la peor forma de conseguirla. Construirla dentro de la
        aplicación operativa la perjudicaría.
\end{itemize}

Se decidirá por separado cuando la aplicación esté en pie. Lo único que este
documento fija sobre ella es que \decision{heredará la paleta oficial}, para que
los dos productos se reconozcan como de la misma casa.

\textbf{Queda fuera del primer alcance, pero decidido}, el comportamiento sin
conexión y la instalación en el dispositivo. No se descartan: la aplicación los
tendrá. Se aplazan, y la sección 7 deja escrito con qué forma entrarán y por qué
aplazarlos no incumple ningún requisito.

\textbf{También queda fuera} todo lo que dependa de módulos de backend que no
existen: reportes de servicio, firma digital, hojas de vida, inventario, alertas
y módulo comercial. Sus pantallas se diseñarán cuando exista contra qué
construirlas.

% =====================================================================
\sectiontitle{2. Lo que la especificación impone}

Esta sección no decide nada. Recoge lo que ya está escrito y que el diseño tiene
que satisfacer, para que las decisiones de las secciones siguientes se lean
contra algo y no en el vacío.

\subsectiontitle{2.1. Los seis requisitos de interfaz}

\begin{longtable}{p{1.8cm} p{6cm} p{6.2cm}}
\textbf{Código} & \textbf{Qué pide} & \textbf{Qué obliga a decidir} \\
\hline
\endhead
RNF-12 & Accesibilidad web & Un nivel concreto; la ERS no fija ninguno \\
RNF-13 & Diseño responsivo & Un orden: móvil primero o escritorio primero \\
RNF-14 & Interacción por pantalla táctil & Un tamaño mínimo de área táctil \\
RNF-15 & Procesos en $\leq$ 6 pasos & Cómo se cuenta un paso \\
RNF-16 & Listas, autocompletado y selección múltiple & De dónde salen esos componentes \\
RNF-17 & Mensajes de error claros & Quién traduce los códigos del backend \\
\end{longtable}

RNF-17 merece una precisión, porque es el único de los seis que tiene ya media
solución construida: el backend emite códigos estables por módulo
(\ruta{ERR_WORK_ORDER_001} y sus hermanos), con la regla explícita de que un
código de ``no existe'' nunca se comparte con uno de ``datos inválidos''. Lo que
falta no es el catálogo sino \textbf{quien lo traduzca a una frase que una
persona lea}. Esa pieza es responsabilidad del frontend y de nadie más.

\subsectiontitle{2.2. Los umbrales de tiempo}

\begin{longtable}{p{3cm} p{10.8cm}}
\textbf{Umbral} & \textbf{Dónde aparece} \\
\hline
\endhead
2 segundos & RNF-03, RNF-04, RNF-09, RNF-10, RNF-18, RNF-22 — respuesta a consultas \\
3 segundos & Apartado 3.3 de la ERS, OBJ-07 y CE-08 del documento de constitución \\
3 minutos & RNF-01 y RNF-20 — procesos completos de usuario \\
\end{longtable}

Los dos primeros son compatibles: dos segundos satisfacen un umbral de tres.
\textbf{Ningún documento lo dice}, y esta lectura es de quien contrasta, no un
hecho declarado. Este documento adopta \decision{el umbral más exigente, dos
segundos}, porque satisface a los dos y evita tener que reconciliarlos.

\subsectiontitle{2.3. El trabajo en campo}

Tres apartados de la ERS describen el escenario real de uso, y juntos son la
restricción de diseño más fuerte del documento:

\begin{itemize}
  \item \textbf{Apartado 1.2}: el sistema ``facilitará el trabajo en campo de
        los ingenieros mediante el uso de dispositivos móviles''.
  \item \textbf{Apartado 2.1}: accesible desde el teléfono del ingeniero
        \textbf{sin instalar aplicaciones nativas}.
  \item \textbf{Apartado 3.1.2}: la pantalla táctil figura entre las interfaces
        de hardware del sistema, junto a la cámara.
\end{itemize}

La consecuencia es que \textbf{el caso principal no es un escritorio}. Un
ingeniero de pie en el pasillo de un hospital, con una mano ocupada, es el
usuario que el diseño debe servir primero.

\subsectiontitle{2.4. Los requisitos que sólo esperan un formulario}

Cuatro requisitos de órdenes de trabajo están abiertos con el backend completo:

\begin{longtable}{p{1.8cm} p{6cm} p{6.2cm}}
\textbf{Código} & \textbf{Qué pide} & \textbf{Qué falta} \\
\hline
\endhead
RF-03 & Formulario de orden: cliente, sede, áreas, tipo y periodicidad & Los cinco datos existen en el API; el formulario no \\
RF-04 & Selección múltiple de áreas de la sede & Es un paso de interfaz. Las áreas elegidas \textbf{no se persisten a propósito}: quedan implícitas en los equipos \\
RF-06 & Visualización de equipos por área & La consulta existe; falta la interfaz que agrupe \\
RF-07 & Selección múltiple de equipos & El API añade de uno en uno; la selección múltiple es del formulario \\
\end{longtable}

Cerrar estos cuatro es el objetivo del primer bloque de construcción (sección 9).

% =====================================================================
\sectiontitle{3. Decisiones de plataforma}

\subsectiontitle{3.1. Framework: Angular}

\decision{La aplicación se construye con Angular.}

\textbf{Por qué.} Angular impone inyección de dependencias y una estructura de
módulos y servicios que \textbf{es el espejo natural de la arquitectura
hexagonal del backend}. Un servicio inyectable que encapsula las llamadas de un
módulo es lo que en el backend es un puerto; un guard de ruta es lo que allí es
una anotación de autorización. Esa correspondencia hace que las dos mitades del
sistema se puedan razonar con el mismo vocabulario, lo que importa en un
proyecto cuya documentación insiste en que las decisiones se comparen entre
capas.

\textbf{Qué cuesta, y queda escrito.} El proyecto original
(\ruta{bolivarbioingenieria-app}) tiene un starter de autenticación en
\ruta{Frontend/src/auth/} —instancia de Keycloak, proveedor de contexto, guard
de ruta y ayudante de peticiones— que la documentación técnica clasificó como
\textbf{directamente reutilizable}. Elegir Angular lo descarta: es React. Hay
que reescribirlo contra \ruta{keycloak-angular}.

Son cuatro piezas pequeñas y el patrón se conserva, pero \textbf{es la primera
vez que este proyecto desecha algo marcado como reutilizable}, y se hace con los
ojos abiertos: la coherencia estructural con el backend se consideró más
valiosa que ahorrar la reescritura de cuatro archivos.

\textbf{Versión.} \decision{No se fija aquí una versión de memoria.} Se adopta
la estable en el momento de crear el proyecto y \textbf{se anota en este mismo
apartado al crearlo}, igual que el backend fija Spring Boot 4.1.1. Una versión
escrita de antemano y nunca comprobada es una afirmación que nace falsa.

\subsectiontitle{3.2. Estilos: Tailwind con la paleta oficial como tokens}

\decision{Tailwind CSS, con los colores oficiales del proyecto declarados como
tokens del tema.}

La paleta no se inventa: ya existe y es la de la ERS, el documento de
constitución y la matriz de trazabilidad. Reutilizarla es lo que hace que la
aplicación y sus documentos se reconozcan como del mismo producto. El detalle
del sistema visual va en la sección 5, incluido \textbf{un conflicto real entre
esa paleta y el nivel de accesibilidad que este documento declara}.

\subsectiontitle{3.3. Componentes: spartan/ui sobre Angular CDK}

\decision{Los componentes de interfaz se toman de spartan/ui.}

El planteamiento importa más que la librería: \textbf{el código de cada
componente se copia al repositorio} en vez de consumirse como caja negra, y
debajo van las primitivas del Angular CDK, que resuelven lo difícil de verdad
—superposiciones, atrapado de foco, navegación de listas por teclado, anuncios
para lectores de pantalla—.

\textbf{Por qué no una librería cerrada.} RNF-14 exige uso táctil y RNF-12
accesibilidad verificable. Las dos cosas obligan a ajustar áreas de pulsación y
comportamiento de foco componente a componente. Con el código en el repositorio
eso es una edición; con una librería cerrada es una pelea contra su sistema de
temas.

\textbf{Qué cubre de RNF-16.} Los tres componentes que el requisito nombra
—listas, autocompletado y selección múltiple— existen como primitivas
accesibles y no hay que escribirlos desde cero.

\textbf{Lo que hay que rehacer, y es trabajo real.} Las plantillas de spartan/ui
traen esquinas redondeadas y una escala de color propias. \textbf{El manual de
marca exige radio cero} y su propia paleta, de modo que cada componente que se
incorpore se re-viste antes de usarse. Que el código viva en el repositorio es
precisamente lo que hace eso posible: con una librería cerrada habría que pelear
contra su sistema de temas para conseguir lo mismo.

\subsectiontitle{3.4. Datos: TanStack Query, aislado tras un servicio por módulo}

\decision{El estado de servidor se gestiona con TanStack Query.}

Resuelve caché, deduplicación de peticiones, reintentos, estados de carga y
error, e invalidación después de escribir. Eso último es lo que impide el
defecto clásico de dos pantallas que muestran la misma lista y sólo una se
entera del cambio.

\textbf{El riesgo, declarado.} El adaptador de Angular se distribuye con el
sufijo \textit{experimental} en el nombre del paquete. Es una dependencia
inestable en el centro de la capa de datos y no se oculta.

\textbf{Cómo se acota.} \decision{Ningún componente usa TanStack Query
directamente.} Cada módulo expone un servicio inyectable propio
—\ruta{WorkOrderApiService} y sus hermanos— y la librería vive dentro de él. Si
su API cambia entre versiones, lo que se toca son esos servicios y no cada
pantalla. Es la misma idea por la que el backend publica puertos en vez de
dejar que un módulo llame al almacén de otro.

\subsectiontitle{3.5. Formularios: los reactivos de Angular}

\decision{Formularios reactivos, con validación declarada en el modelo.}

La razón no es de gusto sino de requisito: \textbf{RNF-07 pide impedir editar
manualmente los datos autocompletados}. Eso no es un atributo visual de
``sólo lectura'' sino una regla del formulario, y un formulario reactivo la
puede expresar y probar sin montar la pantalla.

\subsectiontitle{3.6. Contrato: el cliente se genera desde OpenAPI}

\decision{Los tipos y los clientes de llamada se generan a partir de los
documentos OpenAPI que el backend publica, y no se escriben a mano.}

El backend declara un grupo por módulo y \textbf{tres de ellos ya tienen prueba
que verifica que el grupo contiene sus recursos}. Generar desde ahí extiende esa
garantía al frontend: si el backend renombra un campo, el frontend \textbf{deja
de compilar}. Escribir los tipos a mano convierte ese mismo cambio en un fallo
en tiempo de ejecución, o en ninguno.

\textbf{Lo generado se versiona} en el repositorio. Así el proyecto compila sin
necesidad de un backend levantado, y los cambios del contrato aparecen como
diferencias revisables en lugar de como un resultado distinto del generador.

% =====================================================================
\sectiontitle{4. Arquitectura}

\subsectiontitle{4.1. Organización por módulo de negocio}

\decision{El código se organiza por módulo de negocio, con los mismos nombres
que el backend}, y no por tipo técnico.

\begin{verbatim}
src/app/
  core/         sesion, interceptores, traduccion de errores
  shared/       componentes de interfaz, utilidades
  features/
    client/     equipment/     person/
    location/   workOrder/
\end{verbatim}

\textbf{Por qué no por tipo técnico.} El frontend del proyecto original está
organizado en \ruta{pages/}, \ruta{services/} y \ruta{auth/}, y la documentación
técnica dice explícitamente que \textbf{no hay ahí una convención que copiar} y
que corresponde a MalphasOS decidirla. Repetir esa estructura reparte cada
módulo entre tres carpetas y obliga a saltar entre ellas para leer una función
completa.

\textbf{Qué gana espejar los nombres.} La matriz de trazabilidad relaciona
requisitos con código. Si \ruta{features/workOrder/} se llama igual que el
paquete del backend, la relación se lee sin traducción.

\subsectiontitle{4.2. Capas dentro de un módulo}

\begin{longtable}{p{3.6cm} p{10.2cm}}
\textbf{Capa} & \textbf{Responsabilidad} \\
\hline
\endhead
Cliente generado & Tipos y llamadas, producidos desde OpenAPI. No se edita a mano \\
Servicio del módulo & Encapsula TanStack Query y el cliente generado. Es la única puerta al servidor \\
Componentes & Presentación e interacción. No conocen la librería de datos ni el transporte \\
Modelo de vista & Traduce entre lo que el API devuelve y lo que la pantalla necesita, cuando no coinciden \\
\end{longtable}

La regla que sostiene el conjunto: \decision{un componente nunca llama al API
directamente}. Siempre pasa por el servicio de su módulo.

\subsectiontitle{4.3. Autenticación y autorización}

\textbf{Autenticación} con \ruta{keycloak-angular} sobre el flujo de código de
autorización con PKCE, que es el correcto para un cliente público. El patrón del
proyecto original se conserva aunque su código no: instancia única, inicio de
sesión al arrancar, refresco automático del token antes de que expire y cierre
de sesión si el refresco falla.

\textbf{Autorización.} El backend define 19 autoridades y exige una sola,
nombrada literalmente, en cada una de sus 92 operaciones. El frontend
\decision{refleja esas mismas autoridades en los guards de ruta y en la
visibilidad de los controles}.

Dos precisiones que evitan malentendidos peligrosos:

\begin{itemize}
  \item \textbf{El frontend no autoriza: oculta.} Esconder un botón mejora la
        experiencia; el permiso lo sigue comprobando el servidor en cada
        llamada. Un frontend que ``protege'' una operación no protege nada.
  \item \textbf{El vocabulario no se duplica a mano.} Las autoridades salen del
        token, y la lista se contrasta con la del backend. Dos listas escritas
        por separado se desincronizan, y ésta ya tiene un precedente
        documentado en el propio proyecto.
\end{itemize}

\subsectiontitle{4.4. Traducción de errores}

\decision{Un catálogo por módulo traduce cada código del backend a un mensaje en
español}, con un mensaje genérico de respaldo para los códigos que aún no tenga.

Esto satisface RNF-17 y respeta una distinción que el backend se tomó el trabajo
de mantener: un código de ``no existe'' y uno de ``datos inválidos'' nunca se
comparten, y el de conflicto de estado —arrancar una orden ya ejecutada— es un
tercero distinto de los dos. \textbf{Si el frontend los funde en un ``ha ocurrido
un error'', tira a la basura una decisión de diseño del servidor.}

Las tres familias se presentan de forma distinta: lo que no existe explica qué
referencia falló; lo inválido señala el campo; el conflicto de estado dice qué
impide la operación \textbf{ahora}.

La redacción sigue el tono de voz del manual de marca (5.5): directa e
instructiva, sin disculpas. ``Esta orden ya se ejecutó y no admite cambios'' y
no ``¡Ups! Algo salió mal''.

% =====================================================================
\sectiontitle{5. Sistema visual}

\subsectiontitle{5.1. La fuente es el manual de marca}

\decision{El sistema visual no se define en este documento: lo define el
\textit{Manual de Marca de MalphasOS} v1.0, y aquí sólo se declara cómo se
traduce a una interfaz.}

Vive en \ruta{Documentation/FrontendDesign/BrandManual/}. Es un documento
cerrado y anterior a esta declaración, de modo que \textbf{donde los dos
discrepen, manda el manual}, salvo en lo que la sección 5.3 señala
explícitamente como extensión y no como cambio.

Lo que el manual fija y esta declaración adopta sin discutir:

\begin{itemize}
  \item \textbf{Una paleta monocroma} —tinta y grises sobre papel— con
        \textbf{un solo rojo de acento}, reservado a la acción principal y a
        las alertas, que \textbf{nunca cubre más del 10\,\% de una
        composición}.
  \item \textbf{Radio cero.} Ángulos rectos en todo: botones, campos,
        tarjetas, imágenes.
  \item \textbf{Alineación a la izquierda siempre.} No se centra ni se
        justifica.
  \item \textbf{Una sola familia tipográfica}, Archivo, con tres pesos.
  \item \textbf{Los estados operativos se distinguen por peso tipográfico y
        por regla, no por colores nuevos.}
\end{itemize}

Ese último punto es el que más cambia una interfaz respecto de lo habitual, y
se recoge tal cual: \textbf{una orden vencida no se pinta de rojo}, se compone
con otro peso. El acento señala acciones y alertas, no clasifica estados.

\subsectiontitle{5.2. Color}

\begin{longtable}{p{3.1cm} p{2.1cm} p{8.6cm}}
\textbf{Nombre} & \textbf{Hex} & \textbf{Papel en la interfaz} \\
\hline
\endhead
Tinta & \#201E1D & Texto principal \\
Negro marca & \#2D2B2B & Superficies oscuras, paso 900 de la escala neutra \\
Papel & \#F3F2F2 & Fondo \\
Acento & \#EC3013 & Bordes, anillo de foco, iconografía de alerta, texto grande \\
Acento 700 & \#AE1800 & \textbf{Texto en acento y relleno de la acción principal} \\
\end{longtable}

Más dos escalas de nueve pasos, neutra y de acento, que el manual define y de
las que salen los grises intermedios y los fondos de aviso.

\subsectiontitle{5.2.1. La regla de contraste del manual, y el caso que no nombra}

El manual ya trae su propia regla de contraste: \textit{``El rojo \#EC3013 no se
usa para texto de párrafo sobre papel. Para texto en acento emplear \#AE1800
(paso 700)''}. Medido, funciona:

\begin{longtable}{p{6.6cm} p{2.4cm} p{4.4cm}}
\textbf{Combinación} & \textbf{Contraste} & \textbf{AA texto normal} \\
\hline
\endhead
Tinta \#201E1D sobre papel & 14,86:1 & Cumple \\
Negro marca \#2D2B2B sobre papel & 12,60:1 & Cumple \\
Acento 700 \#AE1800 sobre papel & 6,41:1 & Cumple \\
Acento \#EC3013 sobre papel & 3,76:1 & No cumple — por eso existe la regla \\
\end{longtable}

\textbf{El caso que el manual no cubre es el botón.} Describe la acción
principal como ``relleno acento, etiqueta a la izquierda''. Con el relleno en
\#EC3013 y la etiqueta en papel, el contraste es \textbf{3,76:1}: el control más
repetido de la aplicación no alcanzaría el nivel declarado en la sección 6.

\decision{Resolución: el relleno de la acción principal usa el paso 700,
\#AE1800, con la etiqueta en papel.} Da 6,41:1.

\textbf{No es una corrección del manual sino su propia regla aplicada un caso
más allá}: el manual ya había decidido que el acento puro no sostiene texto
encima; sólo no nombró el relleno sólido entre esos casos. El acento se queda
exactamente donde el manual lo pone para todo lo demás —bordes, anillo de foco,
iconografía de alerta—, y ahí \textbf{sí cumple}: para elementos no textuales el
umbral es 3,0 y \#EC3013 da 3,76:1.

\subsectiontitle{5.3. Tipografía, retícula y espaciado}

Todo lo de esta subsección viene del manual; sólo el último punto es una
extensión propia, y se señala como tal.

\begin{longtable}{p{3.6cm} p{10.2cm}}
\textbf{Aspecto} & \textbf{Regla} \\
\hline
\endhead
Familia & Archivo. Sustitutas: Helvetica Neue, luego Arial \\
Pesos & Títulos 700, subtítulos 600, cuerpo 400 \\
Interlineado & Títulos 0,95--1,05; cuerpo 1,5--1,6 \\
Medida de lectura & Máximo 70 caracteres \\
Retícula & 12 columnas, medianil 24 px, márgenes 48 px \\
Espaciado & Múltiplos de 8 px \\
Radio de esquina & Cero, en todo \\
Iconografía & Lucide, trazo 2 px, esquinas rectas, 16 / 20 / 24 px \\
\end{longtable}

\textbf{Extensión propia:} el manual fija retícula y márgenes pensando en
composiciones de escritorio y no habla de pantallas estrechas.
\decision{En anchos de teléfono la retícula se reduce a 4 columnas y los
márgenes a 16 px, conservando el medianil de 24 px y el espaciado en múltiplos
de 8.} Se declara aquí porque la aplicación es móvil primero (5.4) y el manual
no resuelve ese caso.

\subsectiontitle{5.4. Densidad táctil y orden de diseño}

\begin{itemize}
  \item \decision{Área táctil mínima de 44 por 44 píxeles} para todo control
        accionable. Es lo que convierte RNF-14 en algo medible en lugar de una
        intención. \textbf{Es compatible con el radio cero y con el espaciado de
        8}: 44 no es múltiplo de 8, pero es una superficie mínima de
        interacción, no una medida de maquetación.
  \item \textbf{Móvil primero.} Las pantallas se diseñan para el teléfono y se
        expanden hacia el escritorio. Se sigue del apartado 2.1 de la ERS: el
        teléfono del ingeniero no es un caso degradado, es el caso principal.
  \item \textbf{Tamaño base de 16 píxeles} en el cuerpo. Por debajo, los
        navegadores móviles amplían al enfocar un campo, que es un salto visual
        que confunde.
\end{itemize}

\subsectiontitle{5.5. Tono de los textos de la interfaz}

El manual fija el tono de voz y afecta directamente a RNF-17:
\textbf{directo e instructivo}, nunca entusiasta ni publicitario. Su propio
ejemplo —``Orden cerrada. Próximo mantenimiento: 14 días''— es el modelo para
los mensajes del sistema.

\decision{Los mensajes de error siguen ese tono: dicen qué pasó y qué hacer, sin
disculpas ni signos de exclamación.} La traducción de los códigos del backend
(4.4) se redacta bajo esta regla.

\subsectiontitle{5.6. Cómo se cuenta un paso (RNF-15)}

RNF-15 pide procesos en seis pasos o menos y no define qué es un paso. Este
documento lo define para poder comprobarlo:

\decision{Un paso es cada interacción que el usuario debe completar para
avanzar: abrir una pantalla, rellenar un campo obligatorio, elegir de una lista
o confirmar.} No cuentan la navegación de vuelta, las acciones opcionales ni los
datos que el sistema rellena solo.

Bajo esa definición, crear una orden de trabajo son cinco pasos: cliente, sede,
áreas, equipos, y datos del servicio con la confirmación. Queda uno de margen.

\sectiontitle{6. Accesibilidad}

\decision{La aplicación declara conformidad con WCAG 2.1 nivel AA, y se verifica
automáticamente.}

La ERS pide accesibilidad sin fijar nivel (RNF-12). Sin un nivel, el requisito no
tiene criterio de aceptación y no se puede decir si se cumplió. Este documento lo
fija.

\begin{longtable}{p{5.2cm} p{8.6cm}}
\textbf{Criterio} & \textbf{Cómo se cumple y se comprueba} \\
\hline
\endhead
Contraste $\geq$ 4,5:1 en texto & Acento 700 \#AE1800 para texto y para el relleno de la acción principal (5.2.1) \\
Área táctil $\geq$ 44 px & Tamaño mínimo en los componentes base \\
Foco siempre visible & Indicador de foco propio, nunca suprimido \\
Operable por teclado & Heredado de las primitivas del CDK, verificado en pruebas \\
Formularios etiquetados & Etiqueta asociada y error vinculado a su campo \\
Errores anunciados & Región de anuncio para lectores de pantalla \\
\end{longtable}

\textbf{La verificación no es opcional.} Las pruebas de integración ejecutan un
analizador automático de accesibilidad sobre cada pantalla, y un fallo rompe la
batería. La razón está aprendida en este mismo proyecto: \textbf{una regla
declarada y no ejercitada envejece sola}, y ya ocurrió con una regla de negocio
que se dio por construida durante una jornada entera sin estarlo.

Queda dicho lo que la comprobación automática no cubre: detecta una parte de los
problemas, no todos. El orden de lectura, la claridad del lenguaje y la
utilidad real para un lector de pantalla exigen revisión humana.

% =====================================================================
\sectiontitle{7. Movilidad y comportamiento sin conexión}

\subsectiontitle{7.1. Lo que se construye ahora}

\decision{La aplicación se construye como una aplicación web responsiva y
táctil. No es instalable y no funciona sin conexión.}

\textbf{Las dos capacidades están decididas y aplazadas, no descartadas}: la
aplicación tendrá instalación en el dispositivo y consulta sin conexión para el
trabajo en campo. Se añadirán más adelante, y esta sección deja escrito el
diseño con el que entrarán para que quien lo retome no vuelva a decidirlo.

\subsectiontitle{7.2. Por qué aplazarlo no incumple la especificación}

Conviene que quede dicho, porque el apartado 2.1 de la ERS habla de movilidad y
podría leerse como que exige una aplicación instalable.

No la exige. Lo que pide es que el sistema sea accesible desde el teléfono del
ingeniero \textbf{sin instalar aplicaciones nativas}. \textbf{Una página web
responsiva abierta en el navegador del teléfono satisface ese requisito por
completo}, y de hecho lo satisface de la forma más literal posible: no hay nada
que instalar.

Lo mismo con el apartado 3.1.2, que lista la pantalla táctil entre las
interfaces de hardware: eso lo cubre el diseño táctil de la sección 5.4, que sí
entra desde el primer día.

\textbf{Ningún requisito de la ERS pide instalación ni uso sin conexión.} Son
mejoras del escenario de campo, no obligaciones, y por eso se pueden aplazar sin
dejar nada incumplido.

\subsectiontitle{7.3. Lo que se añadirá después, y con qué forma}

\begin{itemize}
  \item \textbf{Instalación en el dispositivo.} Manifiesto e icono para que el
        ingeniero pueda anclar la aplicación a su pantalla de inicio y abrirla
        a pantalla completa. Sigue sin instalarse nada nativo.
  \item \textbf{Consulta sin conexión, de solo lectura.} Órdenes asignadas al
        ingeniero, con sus equipos, sedes y áreas. Cubre el escenario de campo
        más valioso: \textbf{llegar al hospital sabiendo qué hay que hacer}.
  \item \textbf{Escritura sin conexión: no}, y no por falta de tiempo. Ver 7.4.
\end{itemize}

\subsectiontitle{7.4. Por qué la escritura sin conexión no entrará sin tocar el backend}

Se evaluó permitir crear y avanzar órdenes sin señal, encolando los envíos, y
\textbf{se descartó porque la factura sería del backend}. Lo que sigue está
comprobado sobre el código, no supuesto:

\begin{itemize}
  \item \textbf{No hay claves de idempotencia.} Una orden creada sin señal y
        reintentada al reconectar se crearía dos veces: el identificador lo
        genera hoy el servidor en cada llamada.
  \item \textbf{No hay bloqueo optimista.} No existe columna de versión en
        ninguna de las seis migraciones, de modo que dos ediciones simultáneas
        no se detectan.
  \item \textbf{Reintroduciría una mentira histórica que el diseño ya impide.}
        Una unidad de equipo puede trasladarse de área. Una orden creada sin
        conexión congelaría el área que el dispositivo tenía en caché, que puede
        no ser la real. El módulo de órdenes se diseñó explícitamente para que
        el área congelada sea siempre la verdadera, y la escritura sin conexión
        abriría esa puerta por detrás.
\end{itemize}

\textbf{Mientras la escritura exija conexión, el backend se mantiene cerrado.}
Ésa es la razón de fondo por la que la consulta sin conexión es de solo lectura.

\subsectiontitle{7.5. Decisiones ya tomadas para cuando se retome}

Para que retomarlo sea construir y no volver a decidir:

\begin{itemize}
  \item \textbf{Si algún día entra la escritura sin conexión}, el camino
        decidido es que \decision{el cliente genere el identificador de la orden
        y el backend lo acepte en lugar de generarlo}. Reenviar la misma orden
        escribiría entonces sobre la misma fila. Es un cambio acotado —el
        agregado ya usa identificadores universales— pero toca dominio,
        aplicación y capa de servicios con sus pruebas, y \textbf{el backend
        dejaría de estar cerrado}.
  \item \textbf{Qué se guardará en el dispositivo}: únicamente lo asignado al
        ingeniero que ha iniciado sesión, y se borrará al cerrar sesión. Sin
        cifrado propio: se confía en el del dispositivo, que es lo honesto,
        porque una clave que vive en la misma página que descifra los datos
        protege menos de lo que aparenta.
  \item \textbf{Un supuesto que conviene no perder}: ese límite lo impondría el
        cliente, \textbf{no el servidor}. El backend no filtra por propietario
        —es deuda conocida y declarada— de modo que un usuario con permiso de
        lectura puede solicitar el conjunto completo. La caché respetaría un
        límite que nadie obliga a respetar.
\end{itemize}

\sectiontitle{8. Estrategia de pruebas}

El backend llega a este punto con 624 pruebas, cero fallos y la costumbre de
\textbf{verificar por mutación}: romper la producción a propósito para
comprobar que algo se queja. El frontend adopta el mismo listón.

\begin{longtable}{p{3.4cm} p{10.4cm}}
\textbf{Nivel} & \textbf{Qué cubre} \\
\hline
\endhead
Unitarias & Validación de formularios, guards, traducción de errores, lógica de los servicios de módulo \\
Integración & La pantalla completa contra un API simulado, \textbf{con los códigos de error reales del backend} \\
Extremo a extremo & Inicio de sesión y creación de una orden, contra el sistema real \\
Accesibilidad & Analizador automático sobre cada pantalla, dentro de las pruebas de integración \\
\end{longtable}

Tres decisiones sobre cómo se prueba:

\begin{itemize}
  \item \textbf{El API simulado responde con los códigos reales.} Simular un
        error genérico probaría que la pantalla muestra algo, no que muestra lo
        correcto. La prueba de que \ruta{ERR_WORK_ORDER_002} se convierte en una
        frase útil sólo vale si usa ese código.
  \item \textbf{Las pruebas de extremo a extremo son pocas y deliberadas.} Sólo
        cubren lo que ningún otro nivel puede: que Keycloak y el API se hablen
        de verdad. Es justamente el tipo de integración que compila y no
        funciona.
  \item \textbf{Una prueba que se dobla para pasar describe el defecto en vez
        de detectarlo.} Regla heredada del backend, donde ya evitó cerrar en
        falso un fallo real.
\end{itemize}

% =====================================================================
\sectiontitle{9. Plan de construcción}

\decision{El primer bloque es una rebanada vertical: arranque de la aplicación,
autenticación y el flujo completo de órdenes de trabajo.}

\textbf{Por qué vertical y no por capas.} Atraviesa todo el sistema de una vez
—sesión, datos, formularios, errores, accesibilidad— y cualquier decisión
equivocada de este documento aparece en la primera semana y no en la quinta.

\textbf{Qué arrastra consigo.} El formulario de una orden obliga a construir, de
camino, el selector de cliente, el de sede de ese cliente, la selección múltiple
de áreas de la sede y la de equipos del área. Son piezas que otros módulos
reutilizarán.

\textbf{Qué cierra.} Los cuatro requisitos que hoy están abiertos únicamente por
falta de formulario: RF-03, RF-04, RF-06 y RF-07.

Después, por dependencias: clientes y equipos completos, y luego los módulos que
esperan backend.

% =====================================================================
\sectiontitle{10. Riesgos asumidos y lo que este documento no decide}

\begin{longtable}{p{4.6cm} p{9.2cm}}
\textbf{Riesgo} & \textbf{Por qué se acepta} \\
\hline
\endhead
Dependencia de datos marcada como experimental & Aislada tras un servicio por módulo: si cambia, se tocan los servicios y no las pantallas \\
Se descarta un componente reutilizable & La coherencia estructural con el backend se consideró más valiosa que ahorrar cuatro archivos \\
El acento no alcanza AA como relleno de botón & Medido, no estimado. Resuelto con el paso 700 del propio manual, \#AE1800, que es su regla aplicada a un caso que no nombraba \\
La caché respetaría un límite que el servidor no impone & Declarado en 7.5. No es un riesgo vivo mientras la caché esté aplazada \\
\end{longtable}

\textbf{Lo que este documento deja abierto a propósito:} la versión exacta del
framework, que se fija al crear el proyecto; el diseño de cada pantalla, que
corresponde al momento de construirla; la \textit{landing page}; los módulos
cuyo backend no existe; y \textbf{cuándo entran la instalación y el uso sin
conexión}, que están decididos en forma pero no en fecha.

% =====================================================================
\sectiontitle{11. Trazabilidad}

\begin{longtable}{p{2.4cm} p{11.4cm}}
\textbf{Requisito} & \textbf{Dónde lo atiende este documento} \\
\hline
\endhead
RNF-02 & 3.5 y 3.3 — formularios reactivos y componentes de autocompletado \\
RNF-05 & 9 — una sola entrada de los datos generales en el flujo de la orden \\
RNF-07 & 3.5 — regla del formulario, no atributo visual \\
RNF-12 & 6 — nivel AA declarado y verificado \\
RNF-13 & 5.3 y 5.4 — retícula reducida en teléfono, y móvil primero \\
RNF-14 & 5.4 — área táctil mínima de 44 píxeles \\
RNF-15 & 5.6 — definición de paso y recuento del flujo principal \\
RNF-16 & 3.3 — primitivas accesibles de lista, autocompletado y selección múltiple \\
RNF-17 & 4.4 y 5.5 — catálogo de traducción por módulo, con el tono de voz del manual \\
RNF-03, 04 & 3.4 — caché e invalidación como vía para los umbrales de tiempo \\
RF-03, 04, 06, 07 & 9 — objetivo explícito del primer bloque \\
ERS 2.1 & 7.1 y 7.2 — web responsiva en el navegador del teléfono, sin nada que instalar \\
ERS 3.1.2 & 5.4 — diseño táctil \\
\end{longtable}

\vspace{1em}

Este documento se revisa cuando una de sus decisiones resulte falsa. La
corrección se escribe dejando constancia de qué decía antes y cuándo dejó de ser
cierto, que es la convención del resto de la documentación del proyecto.

\end{document}
